\section{Related Work}
\label{sec:related-work}

This work lies at the intersection of two largely separate lines of research. The first enriches the continuous state that exists between spikes, for example with resonant neurons, oscillatory synaptic kernels, or structured state-space dynamics. The second enriches the timing of discrete events, for example with spike latency, trainable transmission delays, or neuron models whose spike times vary more smoothly with their parameters. The distinction is useful because these mechanisms alter different parts of the hybrid system: oscillatory dynamics change the flow between events, whereas delays change the event-scheduling map. TeS combines one mechanism from each line while retaining a real-valued LIF soma and an event-based exact-gradient formulation.

\paragraph{LIF neurons and exact spike-time gradients.}
For a current-based LIF neuron driven by a first-order exponential synapse, an isolated input produces a difference-of-exponentials postsynaptic-potential kernel, with the alpha function obtained in the equal-time-constant limit. The synaptic and membrane subsystems therefore contribute two real decaying modes. This model is computationally attractive because its state evolution between events is analytic, its threshold condition is simple, and its reset can be represented as a hybrid event. EventProp exploits precisely this structure: it propagates an adjoint through the continuous flow and applies an exact jump condition at each threshold crossing, thereby differentiating spike times without assigning a surrogate derivative to the threshold function \cite{wunderlich2021eventprop}. EventProp changes the learning rule but not the forward neuron dynamics. Its main relevance here is therefore methodological: it establishes that any proposed extension must specify both the augmented continuous flow and the derivatives of every event-time map. At the same time, the standard LIF response remains restricted to a low-pass, single-lobe temporal basis. Weight learning can scale and combine these kernels and can move their threshold crossings, but an individual synapse cannot represent an alternating-sign or phase-selective response.

\paragraph{Resonate-and-fire neurons.}
The resonate-and-fire (RF) neuron replaces the scalar leaky membrane state by a two-dimensional linear system with a complex-conjugate pole pair \cite{izhikevich2001resonate}. Its subthreshold trajectory is a damped rotation, so the response to an input depends on both the elapsed time and the oscillator phase. A sequence of spikes can therefore add constructively or destructively depending on its frequency and relative timing. This supplies band-pass selectivity and a fading phase memory that a one-state LIF neuron does not possess. The change to the forward dynamics is substantial: the somatic state becomes two-dimensional or complex-valued, threshold geometry must be defined in that state space, and reset acts on an oscillator rather than on one leaky voltage. Modern deep RF networks train these models with backpropagation through time and a surrogate derivative for spike generation, rather than with an EventProp-style exact event-time adjoint \cite{alkhamissi2021deeprf,huber2025s5rf}. S5-RF further interprets the RF recurrence as a structured state-space layer and demonstrates that it can be scaled to deep temporal networks \cite{huber2025s5rf}. These results support the computational value of oscillatory poles, but they do not isolate oscillation to the synapse or establish exact continuous-time gradients for a real-soma architecture.

\paragraph{Resonant and oscillatory synaptic kernels.}
Oscillation need not be placed in the spike-generating state. The Synaptic Kernel Inverse Method (SKIM) explicitly used causal temporal basis functions, including damped resonant kernels, to transform incoming spike trains before a thresholded readout \cite{tapson2013skim}. In that construction, the resonant kernels provide a bank of phase-sensitive temporal features, while output weights are obtained by a pseudoinverse or linear regression. The kernel parameters and upstream spike dynamics are not trained end-to-end by a spike-based gradient. More recent work has combined LIF neurons with multiple delayed ``resonant'' synaptic pathways and adapted those pathways to demultiplex inter-spike-interval-coded signals \cite{crook2025synaptic}. There, resonance is implemented as an ISI-selective transmission probability rather than as a complex linear decay; a two-phase local procedure estimates ISI statistics and then aligns synaptic delays. It is not EventProp or surrogate-gradient backpropagation. Relative to a conventional current-based LIF neuron, an oscillatory synapse is a localized forward-model change: the real scalar synaptic trace is replaced by two coupled real traces, while the soma, threshold, and reset may remain unchanged. Taken together, these studies rule out a broad claim that temporally resonant synapses are themselves new. The closer dynamical precedent for TeS is SKIM's damped kernel, while the later work is a precedent for combining synaptic resonance with adaptable delay. TeS differs by placing a complex linear synaptic mode inside a continuously differentiable hybrid SNN and coupling the phase at reception to a state-dependent event delay.

\paragraph{Complex modes in structured state-space models.}
Structured state-space sequence models provide a non-spiking analogue of the same spectral idea. S4 parameterizes a continuous-time linear state-space system whose discretized recurrence or convolution can retain information over long horizons \cite{gu2022s4}. Its state matrix is naturally analyzed through eigenmodes; complex-conjugate modes represent damped oscillations and provide a compact basis for multiple temporal scales and frequencies. Unlike an SNN, an SSM emits dense real-valued activations and is trained by ordinary reverse-mode differentiation through a discretized recurrence or an equivalent convolution. It has neither threshold events nor reset discontinuities. Consequently, the SSM literature does not solve the event-gradient problem, but it gives a strong representation-level motivation for allowing complex poles: a stable complex decay stores both amplitude and phase, whereas an isolated real decay stores only a monotonically fading amplitude. S5-RF makes this connection explicit by deriving a spiking RF architecture from the state-space viewpoint \cite{huber2025s5rf}. TeS borrows the spectral primitive while retaining sparse binary communication and a real threshold variable.

\paragraph{LIF neurons with spike latency.}
The LIF-with-latency (LIFL) model introduces a state-dependent interval between threshold exceedance and spike emission \cite{susi2018lifl}. Below threshold, the internal state behaves as a leaky integrator; above threshold, the neuron enters an active regime and its time-to-fire decreases with the suprathreshold state. Inputs arriving during this active interval can update both the state and the pending spike time. LIFL therefore demonstrates an important precedent: latency can be recomputed from the instantaneous neural state rather than stored as one fixed parameter per axon or synapse. Its forward change is larger than adding a propagation delay to an LIF neuron. Threshold crossing no longer coincides with emission, a second discrete operating mode is introduced, and the neuron remains sensitive to inputs while a spike is pending. The cited LIFL recognition system learns with an online heterosynaptic STDP rule, not with surrogate-gradient backpropagation or an exact spike-time gradient \cite{susi2018lifl}. TeS should consequently not claim state-dependent neural latency as a new concept. Its more specific difference is architectural: emission still occurs at the LIF threshold crossing, the pre-reset slope is sampled at that event, and the resulting state-dependent quantity schedules reception by downstream neurons. Thus the added latency belongs to propagation, not to a suprathreshold spike-generation phase.

\paragraph{Trainable axonal and synaptic delays.}
Trainable delays are an increasingly well-developed route to temporal expressiveness. Discrete-time methods shift spike trains with learnable delay modules or interpolate between time bins, allowing delays to be optimized by backpropagation through time. For example, learnable axonal-delay modules have improved spoken-word recognition using surrogate gradients for threshold events and a finite-difference approximation for the delay derivative \cite{sun2023learnable}. DCLS-based delay learning similarly treats the delayed spike train as a differentiable temporal kernel and trains per-synapse delays with BPTT and surrogate gradients \cite{meszaros2024delaystructure}. In these approaches the LIF flow and reset can remain unchanged; the main forward modification is the time at which a previously emitted event is delivered.

Continuous-time exact-gradient work makes this distinction sharper. DelGrad derives event-based gradients for weights and delays and demonstrates hardware-in-the-loop learning on analog neuromorphic hardware \cite{goeltz2025delgrad}. Delayed EventProp extends exact event-based delay learning to recurrent, multi-spike networks while preserving the memory advantage of event-based adjoints \cite{meszaros2025delay}. These methods are the immediate learning-theoretic baseline for TeS. Their delays are learned model parameters that are fixed for a connection or channel during an inference trajectory; they do not generally vary from one emitted spike to the next as a function of the source neuron's state. TeS instead composes two differentiable timing maps: parameters influence the threshold time and the threshold-crossing slope, and the slope determines a per-event propagation delay. This extra dependency must appear in both the event-time Jacobian and the adjoint jump. The novelty claim should therefore be stated as \emph{state-conditioned exact delay learning}, not merely delay learning.

\paragraph{QIF neurons and loss geometry.}
The quadratic integrate-and-fire (QIF) model changes the membrane flow itself: a quadratic voltage drift replaces the linear leak, and the canonical spike is represented by the finite-time passage of voltage through infinity followed by reinjection, or equivalently by a continuous phase passage in the theta-neuron coordinates. This is a moderate structural change in state dimension---the model still has one somatic degree of freedom---but a major change in the geometry of spike creation. Recent work reports that QIF networks exhibit less fragmented losses and outperform comparable LIF networks under exact spike-based gradient descent \cite{wenig2026qif}. The proposed explanation is directly relevant to TeS: in LIF networks, small parameter changes can alter spike order and trigger spike appearances or disappearances, producing discontinuous loss boundaries and erratic exact gradients; QIF-class dynamics avoid some of these discontinuities. The learning rule in that comparison is an exact spike-based gradient, not a threshold surrogate. QIF therefore provides evidence that improving the forward event geometry can matter as much as improving the differentiation rule. It also creates a necessary caution for TeS. A smooth sigmoid delay map is differentiable within a fixed event topology, but it cannot by itself remove discontinuities caused by changing spike counts or event order. TeS should be evaluated not only for accuracy but also for topology changes, gradient variance, shallow crossings, and delay-induced event reorderings.

\paragraph{Biological relevance of voltage slope.}
There is biological and stochastic-neuron motivation for retaining information about the approach to threshold. In the escape-noise formulation of Plesser and Gerstner, a noisy integrate-and-fire neuron is approximated by a stochastic spike generator whose hazard depends primarily on the distance between the noise-free membrane potential and threshold; the approximation improves when a probability-current term proportional to the voltage derivative is included \cite{plesser2000noise}. In compact form, the instantaneous firing intensity may therefore depend on both $V(t)$ and $\dot V(t)$: at the same subthreshold voltage, a rapidly rising trajectory and a slowly drifting trajectory need not imply the same near-future spike probability. This supports $\dot V$ as a locally meaningful variable for spike timing and provides a biological interpretation for distinguishing sharp from shallow approaches to threshold.

The connection to TeS must nevertheless be stated as inspiration rather than equivalence. The hazard model concerns the probability and expected latency of \emph{spike generation} under noise; it does not show that the conduction time of an already generated action potential is a sigmoid function of the presomatic threshold slope. TeS deterministically transfers this slope information into the propagation stage: sharper crossings are assigned shorter delays, while shallow crossings are assigned longer delays. The mechanism is therefore best presented as a neuromorphic engineering abstraction motivated by slope-sensitive spike generation, not as a direct biophysical model of an axon. A stochastic extension in which the same slope controls an emission or transmission hazard would provide a closer bridge to the biological formulation.

\paragraph{Positioning of TeS.}
The reviewed mechanisms divide naturally into state enrichment and event-time enrichment. RF neurons, resonant synapses, and SSMs add phase-sensitive continuous memory. LIFL, learned delays, and QIF models enlarge or regularize the set of realizable spike times. Prior work establishes each ingredient separately, often with different learning assumptions: SKIM uses linear regression, the cited LIFL system uses local STDP, modern RF and discretized delay networks generally use BPTT with surrogate gradients, and EventProp/DelGrad compute exact event-based gradients. TeS targets the unoccupied intersection: a real LIF soma driven by oscillatory synaptic state, combined with a per-spike delay conditioned on the pre-reset voltage slope, with both mechanisms differentiated in one continuous-time EventProp formulation. Its central computational hypothesis is not merely that oscillation and delay each add capacity, but that the delay controls the phase at which an event enters a resonant filter. This slope-conditioned phase routing is the interaction that the experiments and ablations must isolate.
